\section{模型评价与推广}
\subsection{模型优点}
模型把购电结算、储能动态与信息发布时间统一起来。两日滚动使储能能够跨日转移，避免每晚回到固定电量所产生的人为限制；逐槽因果执行使每一步控制可由当时可得信息实现。
预测权重按历史费用更新，能响应非对称短缺惩罚；候选联合残差块保留负荷、光伏及日内相关结构，并经消融后未进入正式策略。
不同问题共用预测参数、评价区间和初值，有利于比较日内调整与电价信息的影响。

实证结果同时给出主方案、对照、消融与时间留出性能。无对冲更便宜、验证期最优不一定等于开发期最优等负结果被保留，使复杂模块的边际贡献可以检验。
\subsection{局限与改进方向}
首先，单年数据难以覆盖更广的气象与电价状态，100次重采样仅是条件压力评估。应增加独立年份或滚动多折验证，不能据本年度费用承诺未来节省。
其次，评价区间从2月开始，并以6000 kWh重新初始化。进一步研究应从题目给定的1月1日初值连续运行，再提取2月至12月结果，以排除初始化方式的影响。

第三，历史权重成本表仍以单日规划作代理；日内调整也只优化当日剩余时段，并使用零点规划产生的动态日末参考。
将权重标定和日内调整共同迁移到一致的多日滚动目标，可减少层间近似，但需要重新计算正式结果。
两日窗口的最远端尚无剩余电量价值函数，也可能出现窗口末端耗尽倾向；延长窗口或使用经验证的终端价值是值得比较的改进方案。

第四，实际执行器只针对当前缺口充放电，未显式考虑未来价格机会成本。
被删除的场景扩展采用整段情景补救，与在线执行存在信息差异；可进一步研究多阶段非预见性约束或价格感知的短窗实时控制。
最后，当前场景对冲没有表现出冻结验证期或全年降本优势，是否改善新增独立情景下的尾部风险仍需检验；储能衰减和购电容量上限也可在扩展模型中加入。
\subsection{结论}
典型日优化购电费为35126.95元，较无储能节省26.9\%。
在统一334天评价区间内，历史预测的两日滚动计划费用为1395.7万元；
加入组合预报与三点调整后，无对冲主方案为1347.1万元。
波动电价主口径下，两层费用分别为1459.8和1410.3万元。
这些结果支持跨日储能和日内更新的应用价值，同时说明评价复杂调度方法时，应将可执行性、实现费用、尾部风险和信息假设分别检验。
